\section{Apa kesimpulan yang bisa diambil?}\label{sec:discussion}
\subsection{Depresiasi belum menjadi jaminan perbaikan}
Hasil belum memberi dasar untuk mengatakan bahwa rupiah yang melemah selalu memperbaiki trade balance Indonesia. Model total belum mendukung hubungan long run yang kuat. Setelah dipisah, ada pasar atau kelompok yang mengarah pada perbaikan dan ada yang mengarah pada pemburukan.

Pada ARDL, komoditas total dan energi/tambang memberi hasil yang relevan bagi hipotesis awal. Arah negatifnya bertahan pada beberapa pembanding. Namun NARDL belum menguatkan hasil yang sama, dan sebagian pembanding mempunyai catatan diagnostic tests.

Pangsa komoditas yang besar membuat analisis menurut barang berguna. Tetapi pangsa itu belum membuktikan bahwa harga global menyebabkan koefisien kurs tertentu. Model belum memisahkan harga, volume, impor, dan perubahan komposisi.

\subsection{Tiga kesimpulan yang belum bisa dibuat}
Pertama, hasil rasio belum sama dengan hasil ekspor saja. Ekspor bisa turun tetapi rasio membaik jika impor turun lebih cepat. Ekspor bisa naik sementara rasio memburuk jika impor naik lebih cepat.

Kedua, tidak signifikan belum berarti tidak ada efek. Jika confidence interval lebar, arah dan besarnya justru belum diketahui dengan cukup jelas. Untuk menyatakan efek kecil, perlu batas yang masuk akal secara ekonomi dan pengujian khusus.

Ketiga, signifikan pada satu kelompok dan tidak signifikan pada kelompok lain belum membuktikan bahwa keduanya berbeda secara statistik. Selisih koefisien perlu diuji langsung, dengan memperhatikan periode data dan keterkaitan kelompok.

\subsection{Apa keterbatasan yang paling berpengaruh?}
Versi HS per angka belum diketahui. Mapping yang dipakai menghindari pemaksaan kategori, tetapi belum menghapus masalah perubahan kode dan komposisi.

Periode tiap mitra tidak sama. India lebih pendek, sedangkan AS dan Vietnam berhenti lebih awal. Jadi, perbedaan hasil juga mungkin berkaitan dengan kejadian ekonomi yang berbeda dalam sampelnya.

Perdagangan adalah nilai nominal dan belum disesuaikan musiman. Harga bisa mengubah nilainya meskipun volume tidak banyak berubah. Musiman juga masih perlu diperiksa.

Kurs bisa merespons trade balance atau keduanya merespons kejadian yang sama. Lag belum otomatis memisahkan sebab-akibat. Run ini belum memakai desain kausal yang mengisolasi perubahan kurs dari pengaruh lain.

Banyak model juga diuji. BH membantu membaca sebagian pengujian, tetapi belum mengoreksi seluruh pemilihan model, lag, dan sensitivitas. Menampilkan semua hasil membuat proses transparan, bukan berarti semua risiko pencarian hasil hilang.

\subsection{Apa langkah berikutnya?}
Prioritasnya adalah memperbaiki kesebandingan data dan model. Versi HS perlu diperjelas, hasil pada periode yang sama perlu dibandingkan, dan masalah musiman serta residual perlu ditinjau sesuai diagnostic tests.

Jika fokus akhirnya ekspor, persamaan ekspor terpisah akan lebih langsung menjawabnya. Data kuantitas atau harga yang sesuai juga membantu membedakan perubahan nilai dari volume.

Regresi per kode HS-6 memerlukan cek tambahan. Produk bisa mempunyai nol, perubahan kode, atau data terlalu pendek. Data HS-6 yang tersedia belum berarti semua kode bisa langsung memakai model log yang sama.

\subsection{Ringkasan akhir}
Respons perdagangan Indonesia tidak seragam. Australia dan Malaysia menunjukkan hasil NARDL yang lebih mendukung, tetapi ketahanan signifikansinya berbeda. Komoditas total dan energi/tambang mengarah negatif pada ARDL, namun hasil nonlinear dan robustness checks belum seragam.

Pola J terlihat pada sebagian grafik, tetapi confidence interval dan pemeriksaan model mengurangi jumlah kasus yang lebih kuat. Penelitian belum membuktikan bahwa depresiasi selalu merugikan atau komoditas pasti tidak sensitif terhadap kurs.

Saat membaca, pisahkan angka yang terlihat, seberapa jauh model mendukungnya, dan kemungkinan penjelasan ekonominya. Ini membantu kita memakai temuan tanpa memberi kesimpulan yang lebih kuat daripada datanya.
